\begin{warning}[What the consistency theorem does not say]
Neither construction proves \(\mathsf{Con}(\ZFC)\) without a consistency
hypothesis.  The exact second incompleteness result in
\cref{thm:godel-second} applies to a consistent computably enumerable theory
\(T\supseteq\mathsf Q\) and the fixed provability predicate satisfying all
three derivability conditions; under those hypotheses, \(T\) does not derive
its displayed arithmetized consistency sentence.
\end{warning}

\begin{warning}[Countable transitive structures]
The elementary presentation ``take a countable transitive structure
\((M,\in)\) satisfying \(\ZFC\) and construct an \(M\)-generic filter''
assumes more than the bare
arithmetical statement \(\mathsf{Con}(\ZFC)\).  Countability permits an
external enumeration of the dense sets in \(M\), and transitivity makes the
name interpretation transparent.  A proof of the displayed
relative-consistency implication must discharge that pedagogical assumption
by a
syntactic, Boolean-valued, or suitable model-theoretic version of forcing.
The cited relative-consistency theorem for the negation of \(\mathsf{CH}\)
above supplies exactly the Boolean-valued discharge.
\end{warning}

